\begin{ex}
Đoạn thẳng MN dài 18 cm, I là trung điểm của MN. Hỏi đoạn thẳng MI dài bao nhiêu xăng-ti-mét?
\shortans{9}
\loigiai{
$18 : 2 = 9$ cm.
}
\end{ex}

\begin{ex}
Hình tròn có bán kính 6 cm thì đường kính là bao nhiêu xăng-ti-mét?
\shortans{12}
\loigiai{
$6 \times 2 = 12$ cm.
}
\end{ex}

\begin{ex}
Hình tròn có đường kính 24 cm thì bán kính là bao nhiêu xăng-ti-mét?
\shortans{12}
\loigiai{
$24 : 2 = 12$ cm.
}
\end{ex}

\begin{ex}
Một hình vuông có cạnh dài 8 cm. Chu vi của hình vuông đó là bao nhiêu xăng-ti-mét?
\shortans{32}
\loigiai{
$8 \times 4 = 32$ cm.
}
\end{ex}

\begin{ex}
Một hình chữ nhật có chiều dài 14 cm, chiều rộng 6 cm. Chu vi hình chữ nhật đó là bao nhiêu xăng-ti-mét?
\shortans{40}
\loigiai{
$(14 + 6) \times 2 = 40$ cm.
}
\end{ex}

\begin{ex}
Một hình tam giác có độ dài 3 cạnh lần lượt là 7 cm, 9 cm, 11 cm. Chu vi hình tam giác đó là bao nhiêu xăng-ti-mét?
\shortans{27}
\loigiai{
$7 + 9 + 11 = 27$ cm.
}
\end{ex}

\begin{ex}
Khối lập phương có bao nhiêu đỉnh?
\shortans{8}
\loigiai{
Có 8 đỉnh.
}
\end{ex}

\begin{ex}
Khối hộp chữ nhật có bao nhiêu cạnh?
\shortans{12}
\loigiai{
Có 12 cạnh.
}
\end{ex}

\begin{ex}
Một hình vuông có chu vi 28 cm. Cạnh hình vuông đó dài bao nhiêu xăng-ti-mét?
\shortans{7}
\loigiai{
$28 : 4 = 7$ cm.
}
\end{ex}

\begin{ex}
Một hình chữ nhật có chu vi 36 cm, chiều rộng là 8 cm. Chiều dài hình chữ nhật đó là bao nhiêu xăng-ti-mét?
\shortans{10}
\loigiai{
Nửa chu vi = $36 : 2 = 18$ cm. Chiều dài = $18 - 8 = 10$ cm.
}
\end{ex}

\begin{ex}
Một hình chữ nhật có nửa chu vi là 25 cm, chiều dài là 15 cm. Chiều rộng hình chữ nhật đó là bao nhiêu xăng-ti-mét?
\shortans{10}
\loigiai{
$25 - 15 = 10$ cm.
}
\end{ex}

\begin{ex}
Một khối lập phương có cạnh 6 cm. Tổng độ dài các cạnh của khối lập phương đó là bao nhiêu xăng-ti-mét?
\shortans{72}
\loigiai{
$6 \times 12 = 72$ cm.
}
\end{ex}

\begin{ex}
Hình chữ nhật có mấy góc vuông?
\shortans{4}
\loigiai{
Hình chữ nhật có 4 góc vuông.
}
\end{ex}

\begin{ex}
Có bao nhiêu góc vuông trong một hình vuông?
\shortans{4}
\loigiai{
Hình vuông có 4 góc vuông.
}
\end{ex}

\begin{ex}
Đoạn thẳng AB dài 30 cm, điểm M là trung điểm của AB. Đoạn MB dài bao nhiêu xăng-ti-mét?
\shortans{15}
\loigiai{
$30 : 2 = 15$ cm.
}
\end{ex}

\begin{ex}
Compa mở hai đầu 7 cm vẽ một đường tròn. Đường kính đường tròn đó bằng bao nhiêu xăng-ti-mét?
\shortans{14}
\loigiai{
$7 \times 2 = 14$ cm.
}
\end{ex}

\begin{ex}
Một hình tam giác đều có chu vi 24 cm. Mỗi cạnh của hình tam giác đó dài bao nhiêu xăng-ti-mét?
\shortans{8}
\loigiai{
$24 : 3 = 8$ cm.
}
\end{ex}

\begin{ex}
Một hình tứ giác có chu vi 35 cm, độ dài 3 cạnh là 8 cm, 9 cm, 10 cm. Cạnh còn lại dài bao nhiêu xăng-ti-mét?
\shortans{8}
\loigiai{
$35 - (8 + 9 + 10) = 35 - 27 = 8$ cm.
}
\end{ex}

\begin{ex}
Đồng hồ chỉ mấy giờ thì hai kim vuông góc với nhau vào buổi sáng trước 6 giờ?
\shortans{3}
\loigiai{
Lúc 3 giờ đúng.
}
\end{ex}

\begin{ex}
Một khối hộp chữ nhật có bao nhiêu mặt?
\shortans{6}
\loigiai{
Có 6 mặt.
}
\end{ex}

\begin{ex}
Tính: $6 \times 8 + 15 = \dots$
\shortans{63}
\loigiai{
$48 + 15 = 63$.
}
\end{ex}

\begin{ex}
Tính: $7 \times 9 - 25 = \dots$
\shortans{38}
\loigiai{
$63 - 25 = 38$.
}
\end{ex}

\begin{ex}
Tính: $64 : 8 + 42 = \dots$
\shortans{50}
\loigiai{
$8 + 42 = 50$.
}
\end{ex}

\begin{ex}
Tìm $x$: $x \times 7 = 49$.
\shortans{7}
\loigiai{
$x = 49 : 7 = 7$.
}
\end{ex}

\begin{ex}
Tìm $x$: $x : 6 = 9$.
\shortans{54}
\loigiai{
$x = 9 \times 6 = 54$.
}
\end{ex}
